\section{Introduction}\label{sec:intro}

In small-batch large language model (LLM) decoding, projection kernels have
little opportunity to amortize weight reads across tokens. Comparing the
estimated weight traffic with a measured streaming bandwidth ceiling can
therefore give a first estimate of step time
\cite{williams2009roofline,deakin2016babelstream}. Finding the source of the
remaining time requires more context. The same GPU kernel can implement
projections in many layers, and its runtime can depend on the cache state
left by earlier launches. A kernel-name summary combines these instances.
A timeline preserves their order, but identifying the model tensors associated
with each instance requires additional information.

Consider a batch-one decode trace from vLLM on an NVIDIA DGX Spark. A
kernel-name summary puts 148 CUTLASS FP8 launches per step in one row.
This includes seven weight groups and one launch without captured parameters
(\cref{tab:namelevel}). Two groups share the same launch grid and weight
shape. The median runtime was 194.6~\(\mu\)s for Gated DeltaNet (GDN)
output projections and 147.3~\(\mu\)s for attention output projections. Four representative weights from the two groups took
147.5~\(\mu\)s each when tested in isolation. Writing 8~MiB before the
same isolated GEMM increased its time to 187.4~\(\mu\)s
(\cref{sec:split}). In a separate read-versus-write control, L2 data-bank
accesses during the GEMM rose with preceding writes while L2 requests and
fills stayed flat. GB10 does not expose DRAM counters, so this supports
the write-back explanation indirectly (\cref{sec:ncu}). In the engine,
the recurrent update produces several versions of the state for speculative
decoding, leaving dirty lines in L2 that can be evicted as the following
projection streams its weights.

We present \sys{}, a tool and diagnostic workflow for per-launch accounting
in steady-state decoding. \sys{} joins GPU timing records with available
CUDA kernel arguments, including arguments captured from instantiated CUDA
graph nodes. It matches pointer values against an engine-exported memory map
to associate launches with named model weights and cache regions. In the
low-token regime studied here, the size of a referenced weight estimates
the bytes streamed by its projection. Comparing this estimate with the
bandwidth ceiling helps prioritize launches for investigation.
\sys{} also reports repeated references to the same weight and compares
timings across weight groups that share a kernel, grid, shape, and data
type. Nearby launches that reference cache regions provide candidates for
follow-up experiments that test whether preceding work accounts for a
timing difference.

We make three contributions.

\begin{itemize}
  \item A diagnostic workflow that combines a measured bandwidth baseline
    with matched-launch comparisons and controlled changes to the preceding
    workload to test for cross-kernel effects.
  \item An open-source implementation\footnote{\url{https://github.com/sec4si/ledger}}
    that joins CUDA launch timing and argument data with model tensor
    identities and reports timing divergences and repeated weight
    references.
  \item A case study of Qwen3.8-27B NVFP4 with three speculative draft
    tokens on DGX Spark. We first identified the projection anomaly
    manually and subsequently reproduced it with \sys{}. We ported an
    existing deferred-state-store scheme from the Triton path to the
    fused CUDA path and submitted it to vLLM
    \cite{vllmpr61434,vllmpr69160}. The CUDA change preserved
    bitwise equality in kernel tests and greedy output tokens in the
    tested decode runs. It increased serving throughput by 5.2\% and
    7.5\% at concurrency four and eight, respectively, on this system.
\end{itemize}
